%\chapter{Environment Design \& Implementation}\label{sec:designImplementation}


%This chapter should defines where to optimize (the environment, parameters, mutation operators) — \textbf{the domain}


%\textbf{Describe the parameterized cooperative environment(s) $\rightarrow$ Overcooked, \url{https://minigrid.farama.org/environments/babyai/GoTo/},}

%\begin{itemize}
    %\item \textbf{the underspecified dimensions:}
    %\begin{itemize}
        %\item \textbf{Structural parameters:} Grid dimensions, wall placement (which cells are walls vs.\ open), room connectivity (number and position of doorways/passages). Layout encoding — as a 2D grid of tile types, a flattened vector, or an ASCII string.
       % \item \textbf{Functional object placement:} Positions of pots, ingredient dispensers (onion, tomato, broccoli), plate dispensers, serving stations. These determine the task structure and which agent has natural access to which resource.
        %\item \textbf{Agent configuration:} Start positions and orientations of agents. Optionally, view radius (controlling partial observability), allocation radius, field of view.
        %\item \textbf{Task parameters (if applicable):} Recipe set, number of required ingredients per dish, episode length.
   % \end{itemize}

  %  \item \textbf{Level Generation and Mutation Operators}
  %  \begin{itemize}
  %      \item \textbf{(Random generation (for DR and PLR):} How are random valid layouts sampled? Uniform over tile types? Structured generation (place walls with some probability, then place objects)?
  %      \item \textbf{Mutation operators:} What constitutes an ``edit'' to a level? Adding/removing a wall tile, moving a functional object, swapping two tiles, changing agent start positions.
  %      \item \textbf{Difussion}
  %  \end{itemize}

  %  \item \textbf{and the technical infrastructure (likely JAX-based):}
   % \begin{itemize}
   %     \item \textbf{environment implementation:} JaxMARL's Overcooked (or your custom extension), JAX-based, compatible with \texttt{jit}/\texttt{vmap} for vectorized rollouts.
    %    \item \textbf{UED pipeline:} JaxUED for PLR/ACCEL logic — level buffer, regret scoring, replay decisions, mutation.
    %    \item \textbf{Training algorithm:} IPPO or MAPPO via JaxMARL, with parameter sharing.
    %    \item \textbf{Integration:} How these components connect — the outer loop (sample/replay level $\rightarrow$ rollout with IPPO $\rightarrow$ compute regret score $\rightarrow$ update buffer $\rightarrow$ update policy) or similar.
   % \end{itemize}
%\end{itemize}